% Word-Like CV Template
% Centralized style layer. Edit this file to adjust the whole resume.

\usepackage[empty]{fullpage}
\usepackage{fix-cm}
\usepackage{fontawesome5}
\usepackage[dvipsnames]{xcolor}
\usepackage{enumitem}
\usepackage[hidelinks]{hyperref}
\usepackage[normalem]{ulem}
\usepackage{fancyhdr}
\usepackage{tabularx}
\usepackage{graphicx}
\usepackage{array}
\usepackage{ifthen}
\usepackage{fontspec}
\usepackage{xeCJK}

\setmainfont[
  BoldFont = lmroman10-bold.otf,
  ItalicFont = lmroman10-italic.otf,
  BoldItalicFont = lmroman10-bolditalic.otf
]{lmroman10-regular.otf}

\setCJKmainfont[
  BoldFont = FandolSong-Bold.otf,
  ItalicFont = FandolKai-Regular.otf
]{FandolSong-Regular.otf}

\pagestyle{fancy}
\fancyhf{}
\fancyfoot{}
\renewcommand{\headrulewidth}{0pt}
\renewcommand{\footrulewidth}{0pt}
\setlength{\footskip}{8pt}

\addtolength{\oddsidemargin}{-0.6in}
\addtolength{\evensidemargin}{-0.6in}
\addtolength{\textwidth}{1in}
\addtolength{\topmargin}{-.65in}
\addtolength{\textheight}{1.15in}

\urlstyle{same}
\setlength{\tabcolsep}{0pt}
\setlength{\parindent}{0pt}
\setlength{\parskip}{0pt}
\raggedbottom
\raggedright
\sloppy

% -------------------------
% Global style parameters
% -------------------------
\definecolor{CvBlack}{RGB}{0,0,0}
\definecolor{CvPlaceholder}{RGB}{245,245,245}
\definecolor{CvPlaceholderBorder}{RGB}{150,150,150}

\newlength{\cvSectionAfter}
\newlength{\cvSectionBefore}
\newlength{\cvHeadingTwoAfter}
\newlength{\cvHeadingThreeAfter}
\newlength{\cvBodyAfter}
\newlength{\cvBulletAfter}
\newlength{\cvSectionRuleGap}
\newlength{\cvBulletLeft}
\newlength{\cvHeadingRightWidth}

\setlength{\cvSectionAfter}{2.5pt}
\setlength{\cvSectionBefore}{4pt}
\setlength{\cvHeadingTwoAfter}{1.5pt}
\setlength{\cvHeadingThreeAfter}{1.5pt}
\setlength{\cvBodyAfter}{1.5pt}
\setlength{\cvBulletAfter}{1.2pt}
\setlength{\cvSectionRuleGap}{1.5pt}
\setlength{\cvBulletLeft}{1.55em}
\setlength{\cvHeadingRightWidth}{0.30\textwidth}

\newcommand{\cvNameFont}{\fontsize{26pt}{30pt}\selectfont\bfseries}
\newcommand{\cvContactFont}{\fontsize{11pt}{13pt}\selectfont}
\newcommand{\cvSectionFont}{\large}
\newcommand{\cvHeadingTwoFont}{\fontsize{11.5pt}{13.5pt}\selectfont}
\newcommand{\cvHeadingThreeFont}{\normalsize}
\newcommand{\cvBodyFont}{\small}
\newcommand{\cvSmallBodyFont}{\footnotesize}

% -------------------------
% Shared helpers
% -------------------------
\newcommand{\seticon}[1]{\textcolor{CvBlack}{\csname #1\endcsname}}
\newcommand{\cvAfter}[1]{\par\vspace{#1}}
\newcommand{\cvSep}{\hspace{0.7em}|\hspace{0.7em}}

\newcommand{\cvImageOrPlaceholder}[3]{%
  \IfFileExists{#1}{%
    \includegraphics[width=#2]{#1}%
  }{%
    \fcolorbox{CvPlaceholderBorder}{CvPlaceholder}{%
      \parbox[c][#2][c]{#2}{\centering\cvSmallBodyFont #3}%
    }%
  }%
}

% \cvHeader{left image path}{name}{contact line}{right image path}
\newcommand{\cvHeader}[4]{%
  \begin{tabularx}{\textwidth}{@{}m{2.7cm}X m{2.7cm}@{}}
    \centering\cvImageOrPlaceholder{#1}{2.35cm}{Logo}
    &
    \centering
    {\cvNameFont #2}\par
    \vspace{0.85em}
    {\cvContactFont #3}
    &
    \centering\cvImageOrPlaceholder{#4}{2.35cm}{Photo}
  \end{tabularx}
  \cvAfter{8pt}
}

% -------------------------
% Word-like semantic blocks
% -------------------------
\newcommand{\cvSection}[1]{%
  \cvAfter{0pt}%
  {\cvSectionFont #1}\par
  \vspace{\cvSectionRuleGap}%
  {\color{CvBlack}\hrule height 0.45pt}%
  \cvAfter{\cvSectionAfter}%
}

\newcommand{\cvSpacedSection}[1]{%
  \cvAfter{\cvSectionBefore}%
  \cvSection{#1}%
}

\newcommand{\cvHeadingTwo}[2]{%
  {\cvHeadingTwoFont
    \begin{tabular*}{\textwidth}{@{}l@{\extracolsep{\fill}}r@{}}
      #1 & #2 \\
    \end{tabular*}%
  }%
  \cvAfter{\cvHeadingTwoAfter}%
}

\newcommand{\cvHeadingThree}[2]{%
  {\cvHeadingThreeFont
    \ifthenelse{\equal{#2}{}}{%
      \parbox[t]{\textwidth}{#1}%
    }{%
      \begin{tabular*}{\textwidth}{@{}p{\dimexpr\textwidth-\cvHeadingRightWidth-1em\relax}@{\extracolsep{\fill}}r@{}}
        #1 & #2 \\
      \end{tabular*}%
    }%
  }%
  \cvAfter{\cvHeadingThreeAfter}%
}

\newcommand{\cvBodyText}[1]{%
  {\cvBodyFont #1\par}%
  \cvAfter{\cvBodyAfter}%
}

\newcommand{\cvInlineLabel}[2]{%
  \cvBodyText{\textbf{#1}: #2}%
}

\newcommand{\cvInlineLabelZh}[2]{%
  \cvBodyText{\textbf{#1}：#2}%
}

\newcommand{\cvManualBreak}{\par\vspace{\cvBodyAfter}}

\newenvironment{cvBulletList}{%
  \begin{itemize}[
    leftmargin=\cvBulletLeft,
    label=\textbullet,
    labelsep=0.45em,
    itemsep=\cvBulletAfter,
    parsep=0pt,
    topsep=0pt,
    partopsep=0pt
  ]%
  \cvBodyFont
}{%
  \end{itemize}%
  \cvAfter{\cvBodyAfter}%
}

\newcommand{\cvBullet}[1]{\item #1}
\newcommand{\cvPlainBullet}[1]{\item[]\hspace*{-\cvBulletLeft}#1}
